\documentclass[12pt]{exam}
\usepackage[T1]{fontenc}
\usepackage[utf8]{inputenc}
\IfFileExists{lithuanian.ldf}{\usepackage[lithuanian]{babel}}{\usepackage{babel}\babelprovide[main,import]{lithuanian}}
\usepackage{lmodern,graphicx,xcolor,tikz}
\usepackage[a4paper,left=1cm,right=1cm,top=2.5cm,bottom=1cm,footskip=12pt]{geometry}
\usepackage{amsmath,amssymb,tasks}
\newcommand{\irasolaukas}[1]{\begingroup\setlength{\fboxrule}{0.3pt}\setlength{\fboxsep}{3pt}\fcolorbox{black!35}{white}{\rule{0pt}{0.45cm}\hspace{#1}}\endgroup}
\newcommand{\sprendimas}[1]{\par\smallskip\noindent\textit{Sprendimas:}\par\nobreak\vspace{0.1cm}\noindent\begin{tikzpicture}\path[use as bounding box](0,0) rectangle (\linewidth,{-0.7*#1});\foreach \i in {1,...,#1}{\draw[draw=black!25,line width=0.3pt](0,{-0.7*\i}) -- (\linewidth,{-0.7*\i});}\end{tikzpicture}\par\smallskip}
\newcommand{\atsakymas}[1][9]{\par\smallskip\noindent\textit{Atsakymas:}\framebox[#1cm]{\rule{0pt}{0.5cm}}\par}
\pagestyle{headandfoot}\runningheadrule
\pointsdroppedatright\bracketedpoints\setlength{\rightpointsmargin}{1.6cm}
\pointpoints{taškas}{taškai}
\makeatletter
\renewcommand{\points}{%
  \begingroup
    \ifnum\@points=1
      taškas%
    \else
      \ifnum\numexpr\@points-100*(\@points/100)\relax>9
        \ifnum\numexpr\@points-100*(\@points/100)\relax<20
          taškų%
        \else
          \ifnum\numexpr\@points-10*(\@points/10)\relax=1
            taškas%
          \else
            \ifnum\numexpr\@points-10*(\@points/10)\relax>0
              taškai%
            \else
              taškų%
            \fi
          \fi
        \fi
      \else
        \ifnum\numexpr\@points-10*(\@points/10)\relax=1
          taškas%
        \else
          \ifnum\numexpr\@points-10*(\@points/10)\relax>0
            taškai%
          \else
            taškų%
          \fi
        \fi
      \fi
    \fi
  \endgroup
}
\makeatother
\pointname{ \points}
\marginpointname{ \points}
\firstpagefooter{}{\thepage}{45 min. / 30 taškų}
\runningfooter{}{\thepage}{45 min. / 30 taškų}
\setlength{\parindent}{0pt}
\firstpageheader{\small\begin{tabular}{@{}l@{\hspace{0.2cm}}l@{}}Mokinio kodas: & \irasolaukas{2.3cm} \\[0.15cm] Klasė: & \irasolaukas{2.3cm}\end{tabular}}{\small\textbf{Atsiskaitymas 1}\\2026.10.28}{Aristotelio variantas}
\runningheader{\small Kvadratiniai daugianariai ir nelygybės}{}{\small Aristotelio variantas}
\begin{document}
\begin{center}\textbf{Kvadratiniai daugianariai ir nelygybės}\end{center}
Visuose uždaviniuose laikykite, kad $x\in\mathbb R$. Sprendimus pagrįskite.
\medskip
\begin{questions}
\question[2] Mokytojas pateikė kvadratinį daugianarį $P(x)$, kurio diskriminantas neigiamas ir kuris turi minimumą. Be to, žinoma, kad $P(2)=P(8)$.\droppoints
\par\medskip Mokytojas paklausė:\par
\hspace{0.5cm}--- Kam lygu $P(5)$?\par
Stasiukas atsakė:\par
\hspace{0.5cm}--- Minus aštuoniems!\par
Mokytojas pasakė:\par
\hspace{0.5cm}--- Sėsk, du!\par\medskip
Kodėl Stasiukas gavo dvejetą? Atsakymą pagrįskite.
\sprendimas{6}\atsakymas[11]
\vspace{0.3cm}
\question[3] Duotas kvadratinis daugianaris\droppoints
\[P(x)=2(x+4)(x-3)\]
ir tiesinis daugianaris
\[Q(x)=-2x+6.\]
Su kuriomis realiosiomis $x$ reikšmėmis daugianario $P(x)$ reikšmės yra ne didesnės už daugianario $Q(x)$ reikšmes?
\sprendimas{7}\atsakymas[11]
\newpage
\question[10] Duoti du kvadratiniai daugianariai:\droppoints
\[P(x)=- x^{2} + 2 x + 15,\qquad Q(x)=x^{2} - 9.\]
\begin{tasks}[label=\alph*),label-width=1.2em](1)
\task Nustatykite intervalus, kuriuose daugianaris $P(x)$ yra teigiamas, ir intervalus, kuriuose jis yra neigiamas.
\task Nustatykite intervalus, kuriuose daugianaris $Q(x)$ yra teigiamas, ir intervalus, kuriuose jis yra neigiamas.
\task Raskite bendrus intervalus, kuriuose abu daugianariai yra teigiami.
\task Raskite bendrus intervalus, kuriuose abu daugianariai yra neigiami.
\task Raskite visas $x$ reikšmes, su kuriomis abu daugianariai yra neneigiami.
\task Raskite visas $x$ reikšmes, su kuriomis abu daugianariai yra neteigiami.
\task Raskite bendrus intervalus, kuriuose $P(x)>0$ ir $Q(x)<0$.
\task Raskite bendrus intervalus, kuriuose $P(x)<0$ ir $Q(x)>0$.
\end{tasks}

\sprendimas{10}\atsakymas[12]
\question[5] Duotas kvadratinis daugianaris $P(x)=- 2 x^{2} + 12 x - 20$.\droppoints
\par\medskip Dviem skirtingais būdais įrodykite, kad $P(x)<0$ su visomis realiosiomis $x$ reikšmėmis:
\begin{tasks}[label=\alph*),label-width=1.2em](1)
\task išskirdami pilnąjį kvadratą;
\task remdamiesi tik daugianario koeficientais $a$, $b$, $c$ ir diskriminantu.
\end{tasks}
\sprendimas{6}
\vspace{0.35cm}
\question[5] Kvadratinis daugianaris $P(x)=ax^2+bx+c$ nagrinėjamas uždarame intervale $x\in[2;8]$. Žinoma, kad $P(3)=P(7)$, diskriminantas $D>0$, o ekstremumas lygus $9$. Nustatykite visas $x$ reikšmes, su kuriomis šiame intervale $P(x)$ įgyja mažiausią reikšmę, ir visas $x$ reikšmes, su kuriomis įgyja didžiausią.\droppoints
\sprendimas{7}\atsakymas[11]
\question[5] Sukurkite kvadratinį daugianarį $P(x)$, kuris tenkina abi sąlygas:\droppoints
\begin{tasks}[label=\alph*),label-width=1.2em](1)
\task $P(x)>0$ tada ir tik tada, kai $x<-2$ arba $x>4$;
\task daugianario ekstremumas lygus $-27$.
\end{tasks}
Užrašykite gautą daugianarį pavidalu $P(x)=ax^2+bx+c$.
\sprendimas{8}\atsakymas[12]
\end{questions}
\newpage
\begin{center}\large\textbf{Juodraštis}\end{center}
\end{document}
